% Options for packages loaded elsewhere
\PassOptionsToPackage{unicode}{hyperref}
\PassOptionsToPackage{hyphens}{url}
\documentclass[
]{article}
\usepackage{xcolor}
\usepackage[margin=1in]{geometry}
\usepackage{amsmath,amssymb}
\setcounter{secnumdepth}{5}
\usepackage{iftex}
\ifPDFTeX
  \usepackage[T1]{fontenc}
  \usepackage[utf8]{inputenc}
  \usepackage{textcomp} % provide euro and other symbols
\else % if luatex or xetex
  \usepackage{unicode-math} % this also loads fontspec
  \defaultfontfeatures{Scale=MatchLowercase}
  \defaultfontfeatures[\rmfamily]{Ligatures=TeX,Scale=1}
\fi
\usepackage{lmodern}
\ifPDFTeX\else
  % xetex/luatex font selection
\fi
% Use upquote if available, for straight quotes in verbatim environments
\IfFileExists{upquote.sty}{\usepackage{upquote}}{}
\IfFileExists{microtype.sty}{% use microtype if available
  \usepackage[]{microtype}
  \UseMicrotypeSet[protrusion]{basicmath} % disable protrusion for tt fonts
}{}
\makeatletter
\@ifundefined{KOMAClassName}{% if non-KOMA class
  \IfFileExists{parskip.sty}{%
    \usepackage{parskip}
  }{% else
    \setlength{\parindent}{0pt}
    \setlength{\parskip}{6pt plus 2pt minus 1pt}}
}{% if KOMA class
  \KOMAoptions{parskip=half}}
\makeatother
\usepackage{color}
\usepackage{fancyvrb}
\newcommand{\VerbBar}{|}
\newcommand{\VERB}{\Verb[commandchars=\\\{\}]}
\DefineVerbatimEnvironment{Highlighting}{Verbatim}{commandchars=\\\{\}}
% Add ',fontsize=\small' for more characters per line
\usepackage{framed}
\definecolor{shadecolor}{RGB}{248,248,248}
\newenvironment{Shaded}{\begin{snugshade}}{\end{snugshade}}
\newcommand{\AlertTok}[1]{\textcolor[rgb]{0.94,0.16,0.16}{#1}}
\newcommand{\AnnotationTok}[1]{\textcolor[rgb]{0.56,0.35,0.01}{\textbf{\textit{#1}}}}
\newcommand{\AttributeTok}[1]{\textcolor[rgb]{0.13,0.29,0.53}{#1}}
\newcommand{\BaseNTok}[1]{\textcolor[rgb]{0.00,0.00,0.81}{#1}}
\newcommand{\BuiltInTok}[1]{#1}
\newcommand{\CharTok}[1]{\textcolor[rgb]{0.31,0.60,0.02}{#1}}
\newcommand{\CommentTok}[1]{\textcolor[rgb]{0.56,0.35,0.01}{\textit{#1}}}
\newcommand{\CommentVarTok}[1]{\textcolor[rgb]{0.56,0.35,0.01}{\textbf{\textit{#1}}}}
\newcommand{\ConstantTok}[1]{\textcolor[rgb]{0.56,0.35,0.01}{#1}}
\newcommand{\ControlFlowTok}[1]{\textcolor[rgb]{0.13,0.29,0.53}{\textbf{#1}}}
\newcommand{\DataTypeTok}[1]{\textcolor[rgb]{0.13,0.29,0.53}{#1}}
\newcommand{\DecValTok}[1]{\textcolor[rgb]{0.00,0.00,0.81}{#1}}
\newcommand{\DocumentationTok}[1]{\textcolor[rgb]{0.56,0.35,0.01}{\textbf{\textit{#1}}}}
\newcommand{\ErrorTok}[1]{\textcolor[rgb]{0.64,0.00,0.00}{\textbf{#1}}}
\newcommand{\ExtensionTok}[1]{#1}
\newcommand{\FloatTok}[1]{\textcolor[rgb]{0.00,0.00,0.81}{#1}}
\newcommand{\FunctionTok}[1]{\textcolor[rgb]{0.13,0.29,0.53}{\textbf{#1}}}
\newcommand{\ImportTok}[1]{#1}
\newcommand{\InformationTok}[1]{\textcolor[rgb]{0.56,0.35,0.01}{\textbf{\textit{#1}}}}
\newcommand{\KeywordTok}[1]{\textcolor[rgb]{0.13,0.29,0.53}{\textbf{#1}}}
\newcommand{\NormalTok}[1]{#1}
\newcommand{\OperatorTok}[1]{\textcolor[rgb]{0.81,0.36,0.00}{\textbf{#1}}}
\newcommand{\OtherTok}[1]{\textcolor[rgb]{0.56,0.35,0.01}{#1}}
\newcommand{\PreprocessorTok}[1]{\textcolor[rgb]{0.56,0.35,0.01}{\textit{#1}}}
\newcommand{\RegionMarkerTok}[1]{#1}
\newcommand{\SpecialCharTok}[1]{\textcolor[rgb]{0.81,0.36,0.00}{\textbf{#1}}}
\newcommand{\SpecialStringTok}[1]{\textcolor[rgb]{0.31,0.60,0.02}{#1}}
\newcommand{\StringTok}[1]{\textcolor[rgb]{0.31,0.60,0.02}{#1}}
\newcommand{\VariableTok}[1]{\textcolor[rgb]{0.00,0.00,0.00}{#1}}
\newcommand{\VerbatimStringTok}[1]{\textcolor[rgb]{0.31,0.60,0.02}{#1}}
\newcommand{\WarningTok}[1]{\textcolor[rgb]{0.56,0.35,0.01}{\textbf{\textit{#1}}}}
\usepackage{longtable,booktabs,array}
\newcounter{none} % for unnumbered tables
\usepackage{calc} % for calculating minipage widths
% Correct order of tables after \paragraph or \subparagraph
\usepackage{etoolbox}
\makeatletter
\patchcmd\longtable{\par}{\if@noskipsec\mbox{}\fi\par}{}{}
\makeatother
% Allow footnotes in longtable head/foot
\IfFileExists{footnotehyper.sty}{\usepackage{footnotehyper}}{\usepackage{footnote}}
\makesavenoteenv{longtable}
\usepackage{graphicx}
\makeatletter
\newsavebox\pandoc@box
\newcommand*\pandocbounded[1]{% scales image to fit in text height/width
  \sbox\pandoc@box{#1}%
  \Gscale@div\@tempa{\textheight}{\dimexpr\ht\pandoc@box+\dp\pandoc@box\relax}%
  \Gscale@div\@tempb{\linewidth}{\wd\pandoc@box}%
  \ifdim\@tempb\p@<\@tempa\p@\let\@tempa\@tempb\fi% select the smaller of both
  \ifdim\@tempa\p@<\p@\scalebox{\@tempa}{\usebox\pandoc@box}%
  \else\usebox{\pandoc@box}%
  \fi%
}
% Set default figure placement to htbp
\def\fps@figure{htbp}
\makeatother
\setlength{\emergencystretch}{3em} % prevent overfull lines
\providecommand{\tightlist}{%
  \setlength{\itemsep}{0pt}\setlength{\parskip}{0pt}}
\usepackage[]{natbib}
\bibliographystyle{plainnat}
\emergencystretch=2em
\AtBeginDocument{\sloppy}
\usepackage{bookmark}
\IfFileExists{xurl.sty}{\usepackage{xurl}}{} % add URL line breaks if available
\urlstyle{same}
% fallback for those not using the hyperref driver hyperxmp:
\makeatletter
\@ifundefined{xmpquote}{\newcommand{\xmpquote}[1]{#1}}{}
\makeatother
\hypersetup{
  pdftitle={Statistical Inference and Two-Phase Audit Design for Contaminated Measurements in Clinical Research: The catmech R Package},
  pdfkeywords={\xmpquote{clinical audits}, \xmpquote{measurement
error}, \xmpquote{two-phase sampling}, \xmpquote{multiple
imputation}, \xmpquote{survival analysis}, \xmpquote{robust sandwich
covariance}, \xmpquote{R package}, \xmpquote{data provenance}},
  hidelinks,
  pdfcreator={LaTeX via pandoc}}

\title{Statistical Inference and Two-Phase Audit Design for Contaminated
Measurements in Clinical Research: The catmech R Package}
\author{Patrick Iben\\
Independent Biostatistician\\
\href{mailto:patrickiben@gmail.com}{\nolinkurl{patrickiben@gmail.com}}
\textbar{} ORCID: 0009-0005-2443-7973}
\date{November 6, 2026}

\begin{document}
\maketitle
\begin{abstract}
Continuous measurements in multicenter clinical trials, electronic
health record (EHR) disease registries, and observational
epidemiological cohorts are frequently compromised by transcription
mistakes, unit conversion errors, assay calibration shifts, and missing
values. Traditional statistical packages address missing covariates
under Missing at Random (MAR) assumptions, discarding untrusted
observations and imputing from scratch. In source-data verification
(SDV) and clinical audit designs, however, recorded values are present
but unreliable, and a gold-standard reference measurement is audited on
only a fraction of records. Naive analysis of recorded values suffers
attenuation bias, complete-case analysis of audited records inflates
standard errors, and standard multiple imputation needlessly discards
the precision of valid recordings. This paper presents \texttt{catmech},
an R package providing a unified, auditable framework for clinical
measurement verification, rough two-phase audit-size planning,
predictive calibration checking, downstream multi-endpoint regression
pooling, and provenance record generation. We derive the analytical
conditional posteriors for reference-calibrated contamination mixture
models, give a rough, unvalidated audit-size calculation based on
two-phase influence function expansions, and establish cluster-robust
covariance estimators with Barnard-Rubin finite-sample degrees of
freedom for linear models (CV3 jackknife), generalized linear models
(CR1S Huber-White-Liang-Zeger sandwich), and Cox proportional hazards
survival models (event-based model covariance and clustered Lin-Wei
sandwich). We evaluate the package across 10 frozen validation archives
(over 500,000 requests), including the synthetic contamination study
(52,950 requests), two public paired-discrepancy benchmarks on
ClinicalTrials.gov (\(N = 27,208\)) and NHANES 2017--2018
(\(N = 3,908\); 33,800 requests), clustered variance evaluations
(148,500 requests), and facility-level recounts from the Rwanda Health
Management Information System and psychology effect-size recounts
(45,600 requests). Retaining recorded measurements reduces root mean
squared error by 43\% to 48\% over discarding unaudited values on
ClinicalTrials.gov and NHANES. We articulate an evidence-based decision
rule and report all pre-specified negative and inconclusive findings:
the contamination mixture model avoids severe attenuation bias and
narrows confidence intervals by 44\% relative to complete-case analysis
under simulated gross, separable errors, whereas simple linear
imputation is preferred under subtle, diffuse real-world discrepancies.
Finally, \texttt{catmech\_audit\_trail()} records SHA-256 hashes of the
serialized input data and completed datasets, and of the model
specification as canonical JSON, so that a rerun can be compared with
the record; this can support documentation under FDA 21 CFR Part 11 and
ICH E9, but compliance depends on the validated system and procedures
around the analysis.
\end{abstract}

\section{Introduction}\label{introduction}

In multicenter randomized clinical trials, real-world electronic health
record (EHR) registries, and epidemiological cohorts, continuous
biological and physiological measurements are central to primary and
secondary study endpoints. Circulating biomarkers (e.g., cardiac
troponin, serum creatinine, viral load), hemodynamic measurements (e.g.,
systolic blood pressure), and radiographic dimensions (e.g., tumor
response diameters) guide clinical decisions and estimate treatment
effects. However, continuous recordings in multi-site clinical
registries are inevitably subject to data contamination, including assay
lot calibration shifts, clerical transcription errors, decimal slips,
and inconsistent unit conversions
\citep{carroll2006measurement, angelopoulos2023prediction}.

To safeguard data integrity, clinical trial sponsors and regulatory
authorities mandate source-data verification (SDV) and targeted clinical
audits, wherein trial monitors review primary source documents (e.g.,
clinical chart re-abstractions, central laboratory re-assays) for a
subsample of patients \citep{fda2003part11, ich1998e9}. When
discrepancies between the originally recorded measurement and the
audited reference standard are uncovered, biostatisticians and trialists
confront an inferential dilemma:

\begin{enumerate}
\def\labelenumi{\arabic{enumi}.}
\tightlist
\item
  \textbf{Naively trusting recorded measurements}: Treating recorded
  numbers as exact introduces classic errors-in-variables. For
  continuous exposures, non-differential measurement error attenuates
  regression coefficients toward the null, obscuring true treatment
  effects and biological gradients \citep{carroll2006measurement}. For
  non-linear models (e.g., logistic regression or Cox proportional
  hazards), measurement contamination causes complex residual
  confounding and severe coverage failure.
\item
  \textbf{Restricting analysis to audited cases (complete-case
  analysis)}: Analyzing only the audited subset
  (\(\mathcal{A} \subset \{1, \dots, N\}\), typically 10\% to 25\% of
  the cohort) protects against contamination bias but drastically
  reduces effective sample size, inflating standard errors, widening
  confidence intervals, and sacrificing statistical power
  \citep{neyman1938contribution, breslow2003semiparametric}.
\item
  \textbf{Discarding recorded values and applying standard multiple
  imputation}: Imputing unaudited values under standard Missing at
  Random (MAR) assumptions using packages such as \texttt{mice}
  \citep{vanbuuren2011mice, white2011multiple} or \texttt{Amelia}
  \citep{honaker2011amelia} avoids bias under correct specification.
  However, treating recorded values as completely unobserved discards
  the valuable information present in records that were recorded
  accurately \citep{rubin1996multiple}.
\end{enumerate}

To overcome these limitations, we introduce \texttt{catmech}, an R
package for auditable measurement verification, rough two-phase
audit-size planning, predictive calibration diagnostics, and
multi-endpoint model pooling. The package unites two-phase sampling
theory
\citep{cochran1977sampling, breslow2003semiparametric, robins1994estimation}
with reference-calibrated Bayesian mixture modeling
\citep{huang2013simple}, Barnard-Rubin finite-sample inference
\citep{barnard1999small}, and SHA-256 provenance records for
documentation \citep{fda2003part11}.

\section{Statistical theory and
methodology}\label{statistical-theory-and-methodology}

\subsection{The data generating and contamination
framework}\label{the-data-generating-and-contamination-framework}

Let \(i \in \{1, \dots, N\}\) index patients enrolled across \(C\)
clinical sites. Let \(T_i \in \mathbb{R}\) denote the latent true
biological measurement. Conditional on clinical covariates
\(X_i \in \mathbb{R}^p\) (including the analysis outcome and predictors
in the imputation model), the latent measurement follows:
\[T_i \mid X_i \sim \mathcal{N}(X_i \beta, \sigma_T^2)\]

Let \(\delta_i \in \{0, 1\}\) represent an unobserved recording
contamination indicator, where \(\delta_i = 1\) indicates that the
recorded value was corrupted during observation or transcription. The
contamination probability
\(\pi_{s_i} = \mathbb{P}(\delta_i = 1 \mid S_i = s_i)\) may depend on
facility strata or clinical monitoring quality flags \(S_i\)
(\texttt{corruption\_strata}), and the additive error distribution may
vary across error strata \(g_i\) (\texttt{error\_strata}). Writing
\((\pi, \mu_e, \sigma_e^2)\) for a single stratum to simplify notation,
the observed recorded measurement \(W_i\) is generated as:
\[W_i = (1 - \delta_i) T_i + \delta_i (T_i + \epsilon_i), \quad \epsilon_i \sim \mathcal{N}(\mu_e, \sigma_e^2)\]
where \(\epsilon_i\) represents an additive error mechanism with mean
recording bias \(\mu_e\) and error variance \(\sigma_e^2\).

In Phase II of the audit design, an independent gold-standard reference
measurement \(R_i\) is audited on a subset
\(\mathcal{A} \subset \{1, \dots, N\}\) with sample size
\(n_{\text{audit}} = |\mathcal{A}|\):
\[R_i = T_i + \eta_i, \quad \eta_i \sim \mathcal{N}(0, \sigma_{\text{ref}}^2)\]
where \(\sigma_{\text{ref}}\) (\texttt{reference\_sd}, defaulting to
\(0\)) is a user-declared sensitivity constant representing known
reference assay noise, with \(\sigma_{\text{ref}}^2 = 0\) corresponding
to an exact recount.

\subsection{Collapsed slice-within-Gibbs sampling posteriors and linear
reference
imputation}\label{collapsed-slice-within-gibbs-sampling-posteriors-and-linear-reference-imputation}

In \texttt{catmech\_contamination()}, the estimated parameter vector is
\(\Theta = (\beta, \sigma_T^2, \pi, \mu_e, \sigma_e^2)\). For exact
recounts (\(\sigma_{\text{ref}} = 0\)), the target regression parameters
receive the standard noninformative Jeffreys prior
\(p(\beta, \sigma_T^2) \propto 1/\sigma_T^2\) (or a flat prior on
\(\beta\) with a scaled inverse-\(\chi^2\) prior on \(\sigma_T^2\) when
\(\sigma_{\text{ref}} > 0\)). Contamination probabilities receive Beta
priors \(\pi \sim \text{Beta}(a_{\pi}, b_{\pi})\), and recording error
parameters receive a normal-scaled-inverse-\(\chi^2\) prior
\(\mu_e \mid \sigma_e^2 \sim \mathcal{N}(\mu_0, \sigma_e^2 / \kappa_0)\)
and \(\sigma_e^2 \sim \text{Scaled-Inv-}\chi^2(\nu_0, \tau_0^2)\).
During Markov Chain Monte Carlo (MCMC) estimation in R, the latent
indicators and true values \((\delta_i, T_i)\) on unaudited rows are
analytically integrated out when updating \(\text{logit}(\pi)\),
\(\mu_e\), and \(\frac{1}{2}\log(\sigma_e^2)\) via univariate
stepping-out slice sampling (supplemented by a scale-ridge slice step
when fewer than 5 discrepancies are observed), followed by QR-based
conditional draws of \((\beta, \sigma_T^2)\) and exact conditional draws
of \((\delta_i, T_i)\).

Specifically, for unaudited patients (\(i \notin \mathcal{A}\)),
integrating out the unobserved latent value
\(T_i \mid X_i \sim \mathcal{N}(X_i \beta, \sigma_T^2)\) under
\(\delta_i = 1\) yields the convolved marginal predictive distribution
\(W_i \mid \delta_i = 1, X_i, \Theta \sim \mathcal{N}(X_i \beta + \mu_e, \, \sigma_T^2 + \sigma_e^2)\).
The conditional posterior probability of recording corruption
\(p_i = \mathbb{P}(\delta_i = 1 \mid W_i, X_i, \Theta)\) is therefore
evaluated analytically at each Gibbs iteration as:
\[p_i = \frac{\pi \phi(W_i; \, X_i \beta + \mu_e, \, \sigma_T^2 + \sigma_e^2)}{\pi \phi(W_i; \, X_i \beta + \mu_e, \, \sigma_T^2 + \sigma_e^2) + (1 - \pi) \phi(W_i; \, X_i \beta, \, \sigma_T^2)}\]
where \(\phi(w; \, m, v)\) denotes the Gaussian probability density
function evaluated at \(w\) with mean \(m\) and variance \(v\).

Conditional on drawing \(\delta_i = 0\) for an unaudited record,
\(T_i = W_i\) deterministically. Conditional on drawing
\(\delta_i = 1\), \(T_i\) is sampled from the conjugate Gaussian
conditional posterior:
\[T_i \mid \delta_i = 1, W_i, X_i, \Theta \sim \mathcal{N}(\mu_{T \mid W}, \sigma_{T \mid W}^2)\]
\[\mu_{T \mid W} = \sigma_{T \mid W}^2 \left( \frac{X_i \beta}{\sigma_T^2} + \frac{W_i - \mu_e}{\sigma_e^2} \right), \quad \sigma_{T \mid W}^2 = \left( \frac{1}{\sigma_T^2} + \frac{1}{\sigma_e^2} \right)^{-1}\]

For audited patients (\(i \in \mathcal{A}\)), \(R_i\) directly anchors
\(T_i\) (Setting \(T_i = R_i\) when \(\sigma_{\text{ref}} = 0\), or
drawing from the Gaussian precision-weighted combination of \(R_i\),
\(W_i\), and \(X_i \beta\) when \(\sigma_{\text{ref}} > 0\)). Multiple
imputations \(T^{(1)}, \dots, T^{(m)}\) are drawn from post-burn-in
iterations after verifying Gelman-Rubin \(\hat{R}\) and effective sample
size convergence thresholds \citep{gelman1992inference}.

As a parsimonious alternative when discrepancies are diffuse rather than
discrete, \texttt{catmech\_reference()} implements Gaussian reference
imputation accelerated in C++17 via \texttt{RcppArmadillo}. When called
with \texttt{use\_recorded\ =\ TRUE} (referred to in our benchmarks as
Simple Linear MI, or \texttt{cole\_mi}), \texttt{catmech\_reference()}
includes the recorded value \(W_i\) directly as a continuous linear
predictor alongside \(X_i\) in the audited calibration regression
\(T_i \mid W_i, X_i \sim \mathcal{N}(\gamma_0 + \gamma_W W_i + X_i \gamma_X, \sigma_{T \mid W, X}^2)\)
without positing a discrete contamination indicator \(\delta_i\). When
called with \texttt{use\_recorded\ =\ FALSE} (the v0.3 default,
\texttt{reference\_normal}), \texttt{catmech\_reference()} discards
\(W_i\) on unaudited rows and imputes \(T_i\) solely from \(X_i\).

\subsection{Two-phase audit sample size
planning}\label{two-phase-audit-sample-size-planning}

Prior to re-abstracting medical charts or commissioning secondary lab
assays, investigators must budget the audit sample size
\(n_{\text{audit}}\). Using semiparametric influence function theory for
two-phase estimators
\citep{neyman1938contribution, cochran1977sampling, breslow2003semiparametric, angelopoulos2023prediction},
the asymptotic variance of an efficient parameter estimator at audit
sampling fraction \(f = n_{\text{audit}} / N\) satisfies:
\[\text{Var}(\hat{\theta}_f) = \frac{1}{N} \text{Var}(\psi_i) + \frac{1 - f}{f N} \text{Var}(\psi_i - \mathbb{E}[\psi_i \mid G_i])\]
where \(\psi_i\) is the efficient influence function of \(\theta\) and
\(G_i = (W_i, Y_i, X_i)\) denotes the Phase I auxiliary variables
observed on the full cohort (or \(G_{0,i} = (Y_i, X_i)\) when discarding
\(W_i\)).

The precision ratio \(\rho(f)\), defined as the ratio of confidence
interval width at audit fraction \(f\) relative to an ideal census audit
(\(f = 1\)), is:
\[\rho(f) = \frac{\text{Width}(f)}{\text{Width}(1)} = \sqrt{1 + k \frac{1 - f}{f}}, \quad k = \frac{\text{Var}(\psi_i - \mathbb{E}[\psi_i \mid G_i])}{\text{Var}(\psi_i)}\]
Because the noise parameter \(k\) is estimated with uncertainty, the
realized precision ratio \(\hat{\rho}(f)\) is stochastic.
\texttt{catmech\_audit\_power()} is a rough, unvalidated planning aid:
for each candidate audit fraction it reports the share of plausible
values of \(k\) for which the precision ratio meets a target
\(\rho \le \rho_{\text{target}}\),
\[\text{Share}(n_{\text{audit}}) = \mathbb{P}\left(\sqrt{1 + \hat{k} \frac{1 - f}{f}} \le \rho_{\text{target}}\right),\]
computed from log-normal Monte Carlo draws (\texttt{n\_sim\ =\ 2000}) of
\(\hat{k}\), either from assumed error-rate and correlation parameters
or calibrated to the bias-corrected and accelerated (BCa) bootstrap
interval bounds produced by \texttt{catmech\_audit\_plan()}. This share,
labeled \texttt{power} in the output, reflects uncertainty about \(k\);
it is not the power of a hypothesis test. Without a pilot plan, the
default \(k\) treats the unexplained variation as if it sat in erroneous
records that the analysis can identify, which is optimistic, so the
function also reports the audit size that linear imputation needs to
meet the target under the double-sampling regression estimator
\citep{cochran1977sampling} at the same correlation, which is often
several times larger. No study has checked these recommendations. In
pre-specified validation testing, \texttt{catmech\_audit\_plan()}
returned a \texttt{NO-GO} verdict because its median error in predicting
interval width at a 5\% audit fraction was 11.3\% to 16.8\% (exceeding
the pre-specified 10\% tolerance); accordingly,
\texttt{catmech\_audit\_plan()} is documented as an exploratory
calculation rather than a validated sample-size guarantee.

\subsection{Multi-endpoint model pooling with Barnard-Rubin degrees of
freedom}\label{multi-endpoint-model-pooling-with-barnard-rubin-degrees-of-freedom}

Parameter estimates \(\hat{\theta}_j\) and variance-covariance matrices
\(\hat{U}_j\) across \(j \in \{1, \dots, m\}\) completed datasets are
combined using Rubin's rules \citep{rubin1987multiple}:
\[\bar{\theta} = \frac{1}{m} \sum_{j=1}^m \hat{\theta}_j, \quad \bar{U} = \frac{1}{m} \sum_{j=1}^m \hat{U}_j, \quad B = \frac{1}{m - 1} \sum_{j=1}^m (\hat{\theta}_j - \bar{\theta})(\hat{\theta}_j - \bar{\theta})^\top\]
\[T = \bar{U} + \left(1 + \frac{1}{m}\right) B, \quad \hat{\lambda} = \frac{(1 + 1/m) \text{diag}(B)}{\text{diag}(T)}\]

Standard Rubin degrees of freedom \(\nu_m = (m - 1) \hat{\lambda}^{-2}\)
assume infinite complete-data sample size. In clinical audits where
sample sizes, cluster counts, or event counts are constrained, this
assumption leads to anti-conservative confidence intervals.
\texttt{catmech} implements the Barnard \& Rubin (1999) small-sample
adjustment \citep{barnard1999small}:
\[\nu_{obs} = \left(\frac{\nu_{com} + 1}{\nu_{com} + 3}\right) \nu_{com} (1 - \hat{\lambda}), \quad \nu_{BR} = \frac{\nu_m \nu_{obs}}{\nu_m + \nu_{obs}}\]
where \(\nu_{com}\) is the complete-data degrees of freedom. Downstream
pooling covers three core clinical endpoint families:

\begin{enumerate}
\def\labelenumi{\arabic{enumi}.}
\tightlist
\item
  \textbf{Continuous endpoints (\texttt{catmech\_pool\_lm})}: When
  \texttt{cluster\ =\ NULL} (unclustered), fits ordinary least squares
  models with standard covariance and \(\nu_{com} = N - p\). When
  \texttt{cluster} is specified across \(C\) clinical sites, computes
  the closed-form delete-one-cluster jackknife (\texttt{cv3}) covariance
  \(\frac{C - 1}{C} \sum_{c=1}^C (\hat{\beta}_{(-c)} - \hat{\beta})(\hat{\beta}_{(-c)} - \hat{\beta})^\top\)
  on each completed dataset and pools with \(\nu_{com} = C - 1\).
\item
  \textbf{Binary and count endpoints (\texttt{catmech\_pool\_glm})}:
  Fits generalized linear models (logistic and Poisson). When
  \texttt{cluster\ =\ NULL}, uses standard inverse-information
  covariance with \(\nu_{com} = N - p\). When \texttt{cluster} is
  specified across \(C\) clinical sites, sets \(\nu_{com} = C - 1\) and
  computes the CR1S finite-sample-corrected Huber--White--Liang--Zeger
  cluster-robust sandwich covariance \citep{liang1986longitudinal}:
  \[\hat{V}_{\text{sandwich}} = \frac{C}{C - 1}\frac{N - 1}{N - p} \hat{J}^{-1} \left(\sum_{c=1}^C s_c s_c^\top\right) \hat{J}^{-1}\]
  where
  \(s_c = \sum_{i \in \text{cluster } c} \frac{\partial \ell_i}{\partial \theta}\)
  are cluster score residuals and \(\hat{J}^{-1}\) is the unclustered
  GLM covariance matrix. Point estimates and intervals are optionally
  exponentiated to yield odds ratios or incidence rate ratios.
\item
  \textbf{Time-to-event survival endpoints
  (\texttt{catmech\_pool\_coxph})}: Fits Cox proportional hazards models
  via \texttt{survival::coxph()}. When \texttt{cluster\ =\ NULL}, uses
  the model-based inverse partial Fisher information matrix and defines
  complete-data degrees of freedom strictly as \(\nu_{com} = d - p\),
  where \(d\) is the number of observed failure events and \(p\) is the
  parameter count. When \texttt{cluster} is specified across \(C\)
  clinical sites, computes the cluster-summed Lin \& Wei (1989)
  score-residual sandwich covariance \citep{lin1989robust} and sets
  \(\nu_{com} = C - 1\). Hazard ratios
  \(\text{HR} = \exp(\bar{\theta})\) and confidence intervals are
  computed using the Barnard-Rubin \(\nu_{BR}\) \(t\)-distribution.
\end{enumerate}

\subsection{Out-of-sample predictive calibration
checking}\label{out-of-sample-predictive-calibration-checking}

To screen for predictive miscalibration,
\texttt{catmech\_holdout\_check()} partitions audited records into \(K\)
stratified (or cluster-grouped) cross-validation folds and refits the
imputation workflow on the remaining \(K - 1\) folds
\(\mathcal{D}_{-k}\). Because the mixture predictive distribution has a
discrete point mass at \(W_i\) (retained with probability \(1 - p_i\))
and uses \(m\) finite completions, \texttt{catmech\_holdout\_check()}
computes the tie-randomized empirical rank Probability Integral
Transform (PIT) \citep{dawid1984present} for each held-out reference
record \(i \in \mathcal{A}_k\):
\[u_i = \frac{\sum_{j=1}^m \mathbb{I}(T_{i,-k}^{(j)} < R_i^*) + U_i \left(\sum_{j=1}^m \mathbb{I}(T_{i,-k}^{(j)} = R_i^*) + 1\right)}{m + 1}, \quad U_i \sim \text{Uniform}(0, 1)\]
where \(R_i^* = W_i\) when \(|W_i - R_i|\) lies within
\texttt{agreement\_tol} under exact recounts
(\(\sigma_{\text{ref}} = 0\)) and \(R_i^* = R_i\) otherwise. The
diagnostic flags \texttt{miscalibration\_detected} if a two-sided
Kolmogorov--Smirnov test of \(u_i\) against \(\text{Uniform}(0, 1)\)
yields \(p < 0.01\), if the 95\% Wilson upper bound on empirical
predictive interval coverage falls below the nominal level, or if the
corruption probability calibration \(z\)-statistic exceeds
\(z_{0.995}\). As shown in our empirical validation below, held-out row
checks have low statistical power to detect regression slope bias, so
\texttt{catmech\_holdout\_check()} is a standalone diagnostic screen
rather than a proof of model validity.

\section{Software architecture and comparison with existing
packages}\label{software-architecture-and-comparison-with-existing-packages}

The \texttt{catmech} package is structured into modular S3 components
designed for computational efficiency, numerical stability, and
auditability:

{\def\LTcaptype{none} % do not increment counter
\begin{longtable}[]{@{}
  >{\raggedright\arraybackslash}p{(\linewidth - 4\tabcolsep) * \real{0.1754}}
  >{\raggedright\arraybackslash}p{(\linewidth - 4\tabcolsep) * \real{0.3509}}
  >{\raggedright\arraybackslash}p{(\linewidth - 4\tabcolsep) * \real{0.4737}}@{}}
\toprule\noalign{}
\begin{minipage}[b]{\linewidth}\raggedright
Module
\end{minipage} & \begin{minipage}[b]{\linewidth}\raggedright
Core functions
\end{minipage} & \begin{minipage}[b]{\linewidth}\raggedright
Mathematical and computational role
\end{minipage} \\
\midrule\noalign{}
\endhead
\bottomrule\noalign{}
\endlastfoot
\textbf{Audit planning} & \texttt{catmech\_audit\_plan()},
\texttt{catmech\_audit\_power()} & Two-phase influence function ratio
estimation (exploratory) and rough, unvalidated audit-size curves. \\
\textbf{Data diagnostics} & \texttt{catmech\_recording\_audit()},
\texttt{catmech\_holdout\_check()} & Stratum-specific discrepancy
auditing (Jeffreys beta intervals) and out-of-sample rank-PIT checks. \\
\textbf{Imputation engines} & \texttt{catmech\_contamination()},
\texttt{catmech\_reference()} & Reference-calibrated collapsed
slice-within-Gibbs mixture model and Gaussian reference MI
(\texttt{RcppArmadillo}). \\
\textbf{Model pooling} & \texttt{catmech\_pool\_lm()},
\texttt{catmech\_pool\_glm()}, \texttt{catmech\_pool\_coxph()},
\texttt{catmech\_compare\_strategies()} & Multi-endpoint pooling with
Rubin rules, Barnard-Rubin df, cluster-robust covariance, and strategy
comparison. \\
\textbf{Provenance} & \texttt{catmech\_sensitivity()},
\texttt{catmech\_audit\_trail()} & Multi-scenario sensitivity envelopes
and SHA-256 provenance records. \\
\end{longtable}
}

While several mature R packages address missing data or measurement
error, none combine two-phase audit-size planning, reference-calibrated
contamination mixture modeling, multi-endpoint pooling with
cluster-robust covariance and Barnard-Rubin degrees of freedom, and
SHA-256 provenance records:

{\def\LTcaptype{none} % do not increment counter
\begin{longtable}[]{@{}
  >{\raggedright\arraybackslash}p{(\linewidth - 10\tabcolsep) * \real{0.2273}}
  >{\centering\arraybackslash}p{(\linewidth - 10\tabcolsep) * \real{0.1727}}
  >{\centering\arraybackslash}p{(\linewidth - 10\tabcolsep) * \real{0.1455}}
  >{\centering\arraybackslash}p{(\linewidth - 10\tabcolsep) * \real{0.1455}}
  >{\centering\arraybackslash}p{(\linewidth - 10\tabcolsep) * \real{0.1455}}
  >{\centering\arraybackslash}p{(\linewidth - 10\tabcolsep) * \real{0.1636}}@{}}
\toprule\noalign{}
\begin{minipage}[b]{\linewidth}\raggedright
Capability
\end{minipage} & \begin{minipage}[b]{\linewidth}\centering
\texttt{catmech}
\end{minipage} & \begin{minipage}[b]{\linewidth}\centering
\texttt{mice} \citep{vanbuuren2011mice}
\end{minipage} & \begin{minipage}[b]{\linewidth}\centering
\texttt{simex} \citep{leeder2008simex}
\end{minipage} & \begin{minipage}[b]{\linewidth}\centering
\texttt{mecor} \citep{nab2021mecor}
\end{minipage} & \begin{minipage}[b]{\linewidth}\centering
\texttt{Amelia} \citep{honaker2011amelia}
\end{minipage} \\
\midrule\noalign{}
\endhead
\bottomrule\noalign{}
\endlastfoot
\textbf{Contamination Mixture} & \textbf{Yes} (Gibbs) & No & No & No &
No \\
\textbf{Gold-Standard Audits} & \textbf{Yes} (Two-phase) & Partial (MAR)
& No (Replicates) & Yes (Reg. Cal.) & Partial (Priors) \\
\textbf{Audit-Size Planning} & \textbf{Yes} (MC, unvalidated) & No & No
& No & No \\
\textbf{Calibration Screen} & \textbf{Yes} (Rank PIT) & No & No & No &
Partial (\texttt{overimpute}) \\
\textbf{Linear Model Pooling} & \textbf{Yes} (BR + CV3) & Yes (BR df) &
Yes (SIMEX) & Yes (Reg. Cal.) & Yes (Rubin df) \\
\textbf{GLM (Logit/Poisson)} & \textbf{Yes} (Sandwich) & Yes (BR df) &
Yes (SIMEX) & No & Yes (Rubin df) \\
\textbf{Cox PH Pooling} & \textbf{Yes} (Lin-Wei) & Partial
(\texttt{coxph}) & No & No & No \\
\textbf{SHA-256 Audit Trail} & \textbf{Yes} & No & No & No & No \\
\end{longtable}
}

In confirmatory clinical trials, primary estimands are frequently
time-to-event survival endpoints or binary response endpoints rather
than continuous linear outcomes. All \texttt{catmech} pooling functions
export \texttt{broom}-style S3 methods (\texttt{tidy()},
\texttt{glance()}) and S3 forest-plot visualization (\texttt{plot()}),
while \texttt{catmech\_compare\_strategies()} compares four benchmark
strategies (\texttt{validation\_subsample} complete-case,
\texttt{naive\_trust\_recorded} naive analysis,
\texttt{linear\_mi\_recorded} using
\texttt{catmech\_reference(use\_recorded\ =\ TRUE)}, and
\texttt{contamination\_model} using \texttt{catmech\_contamination()})
side by side for continuous (\texttt{lm}) and generalized linear
(\texttt{glm}) endpoints.

\section{Illustrative clinical trial
workflow}\label{illustrative-clinical-trial-workflow}

The following reproducible R workflow illustrates the \texttt{catmech}
pipeline across all eight analytical stages on a simulated 10-site
clinical cohort (\(N = 1,000\), 25\% reference audit
\(n_{\text{audit}} = 250\)). Because the lab value is a predictor in
every downstream analysis, the imputation model in stage 4 includes each
analysis outcome: the continuous outcome, the binary adverse event, and,
for the Cox model, the event indicator and the Nelson-Aalen cumulative
hazard \citep{white2009imputing}. Leaving an outcome out biases the lab
value's coefficient toward zero, and the pooling functions warn when one
is missing:

\begin{Shaded}
\begin{Highlighting}[]
\FunctionTok{library}\NormalTok{(catmech)}
\FunctionTok{library}\NormalTok{(survival)}

\CommentTok{\# 1. Simulate a 10{-}site cohort with recording contamination and a 25\% reference audit}
\FunctionTok{set.seed}\NormalTok{(}\DecValTok{42}\NormalTok{)}
\NormalTok{cohort }\OtherTok{\textless{}{-}} \FunctionTok{simulate\_catmech\_lab}\NormalTok{(}\AttributeTok{n =} \DecValTok{1000}\NormalTok{, }\AttributeTok{seed =} \DecValTok{42}\NormalTok{)}
\NormalTok{cohort}\SpecialCharTok{$}\NormalTok{site }\OtherTok{\textless{}{-}} \FunctionTok{rep}\NormalTok{(}\FunctionTok{seq\_len}\NormalTok{(}\DecValTok{10}\NormalTok{), }\AttributeTok{length.out =} \FunctionTok{nrow}\NormalTok{(cohort))}
\NormalTok{audited }\OtherTok{\textless{}{-}} \FunctionTok{sample.int}\NormalTok{(}\FunctionTok{nrow}\NormalTok{(cohort), }\AttributeTok{size =} \DecValTok{250}\NormalTok{)}
\NormalTok{cohort}\SpecialCharTok{$}\NormalTok{lab\_ref }\OtherTok{\textless{}{-}} \FunctionTok{ifelse}\NormalTok{(}\FunctionTok{seq\_len}\NormalTok{(}\FunctionTok{nrow}\NormalTok{(cohort)) }\SpecialCharTok{\%in\%}\NormalTok{ audited, cohort}\SpecialCharTok{$}\NormalTok{lab\_true, }\ConstantTok{NA\_real\_}\NormalTok{)}
\NormalTok{cohort}\SpecialCharTok{$}\NormalTok{adverse\_event }\OtherTok{\textless{}{-}} \FunctionTok{as.integer}\NormalTok{(cohort}\SpecialCharTok{$}\NormalTok{outcome }\SpecialCharTok{\textgreater{}} \FunctionTok{median}\NormalTok{(cohort}\SpecialCharTok{$}\NormalTok{outcome))}
\NormalTok{cohort}\SpecialCharTok{$}\NormalTok{time }\OtherTok{\textless{}{-}} \FunctionTok{rexp}\NormalTok{(}\FunctionTok{nrow}\NormalTok{(cohort), }\AttributeTok{rate =} \FunctionTok{exp}\NormalTok{(}\FloatTok{0.25} \SpecialCharTok{*}\NormalTok{ (cohort}\SpecialCharTok{$}\NormalTok{lab\_true }\SpecialCharTok{{-}} \DecValTok{6}\NormalTok{)))}
\NormalTok{cohort}\SpecialCharTok{$}\NormalTok{event }\OtherTok{\textless{}{-}} \FunctionTok{rbinom}\NormalTok{(}\FunctionTok{nrow}\NormalTok{(cohort), }\DecValTok{1}\NormalTok{, }\FloatTok{0.75}\NormalTok{)}
\NormalTok{cohort}\SpecialCharTok{$}\NormalTok{hazard }\OtherTok{\textless{}{-}}\NormalTok{ mice}\SpecialCharTok{::}\FunctionTok{nelsonaalen}\NormalTok{(cohort, time, event)}

\CommentTok{\# 2. Rough audit{-}size planning aid (not validated)}
\NormalTok{power\_plan }\OtherTok{\textless{}{-}} \FunctionTok{catmech\_audit\_power}\NormalTok{(}
  \AttributeTok{N =} \DecValTok{1000}\NormalTok{, }\AttributeTok{target\_ratio =} \FloatTok{1.25}\NormalTok{, }\AttributeTok{power\_target =} \FloatTok{0.80}\NormalTok{,}
  \AttributeTok{error\_rate =} \FloatTok{0.20}\NormalTok{, }\AttributeTok{correlation =} \FloatTok{0.85}\NormalTok{, }\AttributeTok{fractions =} \FunctionTok{seq}\NormalTok{(}\FloatTok{0.10}\NormalTok{, }\FloatTok{0.50}\NormalTok{, }\AttributeTok{by =} \FloatTok{0.05}\NormalTok{)}
\NormalTok{)}

\CommentTok{\# 3. Discrepancy screening across clinical sites (exact 95\% Jeffreys beta intervals)}
\NormalTok{audit\_screen }\OtherTok{\textless{}{-}} \FunctionTok{catmech\_recording\_audit}\NormalTok{(}
  \AttributeTok{data =}\NormalTok{ cohort, }\AttributeTok{value =} \StringTok{"lab\_value"}\NormalTok{, }\AttributeTok{reference =} \StringTok{"lab\_ref"}\NormalTok{,}
  \AttributeTok{by =} \StringTok{"site"}\NormalTok{, }\AttributeTok{agreement\_tol =} \FloatTok{0.5}
\NormalTok{)}

\CommentTok{\# 4. Reference{-}calibrated contamination mixture model (or catmech\_reference(use\_recorded = TRUE))}
\NormalTok{fit }\OtherTok{\textless{}{-}} \FunctionTok{catmech\_contamination}\NormalTok{(}
  \AttributeTok{data =}\NormalTok{ cohort, }\AttributeTok{value =} \StringTok{"lab\_value"}\NormalTok{, }\AttributeTok{reference =} \StringTok{"lab\_ref"}\NormalTok{, }\AttributeTok{m =} \DecValTok{10}\NormalTok{, }\AttributeTok{seed =} \DecValTok{2026}\NormalTok{,}
  \AttributeTok{predictors =} \FunctionTok{c}\NormalTok{(}\StringTok{"age"}\NormalTok{, }\StringTok{"sex"}\NormalTok{, }\StringTok{"inflammation"}\NormalTok{, }\StringTok{"outcome"}\NormalTok{, }\StringTok{"adverse\_event"}\NormalTok{, }\StringTok{"event"}\NormalTok{,}
                 \StringTok{"hazard"}\NormalTok{), }\AttributeTok{corruption\_strata =} \StringTok{"outlier\_flag"}
\NormalTok{)}

\CommentTok{\# 5. Out{-}of{-}sample predictive calibration check (5{-}fold cross{-}validated rank PIT)}
\NormalTok{calib\_check }\OtherTok{\textless{}{-}} \FunctionTok{catmech\_holdout\_check}\NormalTok{(}\AttributeTok{fit =}\NormalTok{ fit, }\AttributeTok{folds =} \DecValTok{5}\DataTypeTok{L}\NormalTok{, }\AttributeTok{m =} \DecValTok{20}\DataTypeTok{L}\NormalTok{, }\AttributeTok{seed =} \DecValTok{123}\DataTypeTok{L}\NormalTok{)}

\CommentTok{\# 6. Multi{-}endpoint downstream regression pooling (LM with CV3, GLM with CR1S, Cox PH with Lin{-}Wei)}
\NormalTok{pool\_continuous }\OtherTok{\textless{}{-}} \FunctionTok{catmech\_pool\_lm}\NormalTok{(}
  \AttributeTok{result =}\NormalTok{ fit, }\AttributeTok{formula =}\NormalTok{ outcome }\SpecialCharTok{\textasciitilde{}}\NormalTok{ lab\_value }\SpecialCharTok{+}\NormalTok{ age }\SpecialCharTok{+}\NormalTok{ sex, }\AttributeTok{cluster =} \StringTok{"site"}
\NormalTok{)}
\NormalTok{pool\_binary }\OtherTok{\textless{}{-}} \FunctionTok{catmech\_pool\_glm}\NormalTok{(}
  \AttributeTok{result =}\NormalTok{ fit, }\AttributeTok{formula =}\NormalTok{ adverse\_event }\SpecialCharTok{\textasciitilde{}}\NormalTok{ lab\_value }\SpecialCharTok{+}\NormalTok{ age }\SpecialCharTok{+}\NormalTok{ sex,}
  \AttributeTok{family =} \FunctionTok{binomial}\NormalTok{(}\AttributeTok{link =} \StringTok{"logit"}\NormalTok{), }\AttributeTok{cluster =} \StringTok{"site"}\NormalTok{, }\AttributeTok{exponentiate =} \ConstantTok{TRUE}
\NormalTok{)}
\NormalTok{pool\_survival }\OtherTok{\textless{}{-}} \FunctionTok{catmech\_pool\_coxph}\NormalTok{(}
  \AttributeTok{result =}\NormalTok{ fit, }\AttributeTok{formula =} \FunctionTok{Surv}\NormalTok{(time, event) }\SpecialCharTok{\textasciitilde{}}\NormalTok{ lab\_value }\SpecialCharTok{+}\NormalTok{ age }\SpecialCharTok{+}\NormalTok{ sex,}
  \AttributeTok{cluster =} \StringTok{"site"}\NormalTok{, }\AttributeTok{exponentiate =} \ConstantTok{TRUE}
\NormalTok{)}

\CommentTok{\# 7. Side{-}by{-}side comparison of all four audit strategies for continuous or GLM endpoints}
\NormalTok{comp }\OtherTok{\textless{}{-}} \FunctionTok{catmech\_compare\_strategies}\NormalTok{(}
  \AttributeTok{data =}\NormalTok{ cohort, }\AttributeTok{value =} \StringTok{"lab\_value"}\NormalTok{, }\AttributeTok{reference =} \StringTok{"lab\_ref"}\NormalTok{,}
  \AttributeTok{formula =}\NormalTok{ outcome }\SpecialCharTok{\textasciitilde{}}\NormalTok{ lab\_value }\SpecialCharTok{+}\NormalTok{ age }\SpecialCharTok{+}\NormalTok{ sex, }\AttributeTok{m =} \DecValTok{10}\NormalTok{, }\AttributeTok{seed =} \DecValTok{2026}
\NormalTok{)}

\CommentTok{\# 8. Provenance record: SHA{-}256 hashes of the data, specification and completions}
\NormalTok{report }\OtherTok{\textless{}{-}} \FunctionTok{catmech\_audit\_trail}\NormalTok{(}
  \AttributeTok{object =}\NormalTok{ fit, }\AttributeTok{pooled =}\NormalTok{ pool\_survival, }\AttributeTok{format =} \StringTok{"markdown"}\NormalTok{, }\AttributeTok{path =} \StringTok{"audit\_trail.md"}
\NormalTok{)}
\end{Highlighting}
\end{Shaded}

Printing \texttt{pool\_survival} reports pooled log-hazard point
estimates, clustered Lin-Wei sandwich standard errors, Barnard-Rubin
degrees of freedom (with \(\nu_{com} = C - 1 = 9\) for \(C = 10\)
sites), exponentiated hazard ratios (\texttt{estimate\_exp}), and 95\%
confidence intervals:

\begin{Shaded}
\begin{Highlighting}[]
\NormalTok{{-}{-} Catmech Pooled Cox PH Model {-}{-}}
\NormalTok{Completions (m): 10 | Observations (n): 1000 | Events: 735}

\NormalTok{      term estimate std.error     df   lower  upper estimate\_exp lower\_exp upper\_exp}
\NormalTok{ lab\_value   0.1877    0.0386 6.0881  0.0936 0.2818       1.2065    1.0982    1.3255}
\NormalTok{       age   0.0046    0.0042 7.2302 {-}0.0051 0.0144       1.0047    0.9949    1.0145}
\NormalTok{      sexM  {-}0.0574    0.0443 7.3102 {-}0.1614 0.0465       0.9442    0.8510    1.0476}
\end{Highlighting}
\end{Shaded}

\section{Synthetic simulation studies and
benchmarks}\label{synthetic-simulation-studies-and-benchmarks}

To evaluate the statistical properties of \texttt{catmech}, we conducted
pre-specified Monte Carlo simulation studies across 10 validation
archives (over 500,000 total requests). The primary reference-calibrated
contamination study (\texttt{validation/contamination-study}) executed
52,950 planned requests across 183 method-scenario cells spanning 4
primary cells (\texttt{clean}, \texttt{flagged\_errors},
\texttt{missed\_flags}, and \texttt{site\_shift\_declared}, 1,000
cohorts each) and 23 stress cells (150 cohorts each), with cohort size
\(N = 600\), linear regression target slope \(\beta = 0.800\) on the
latent covariate \(T_i\), a 25\% reference audit
(\(n_{\text{audit}} = 150\)), and \(m = 20\) imputations. Seven core
strategies were evaluated across the 27 scenarios (with 3
noisy-reference/assay-noise stress scenarios running only
\texttt{contamination}, and 6 scenarios additionally evaluating standard
and bootstrap-donor predictive mean matching \texttt{reference\_pmm} and
\texttt{reference\_pmm\_abb} to total 52,950 requests), yielding 52,900
returned intervals, 16 expected design refusals (14 in
\texttt{reference\_05pct} and 2 in \texttt{reference\_10pct} due to
sparse strata), 33 MCMC convergence refusals, and 1 singular raking fit.

\subsection{Scenario 1: Unflagged, gross, and misspecified recording
contamination}\label{scenario-1-unflagged-gross-and-misspecified-recording-contamination}

The table below compares the seven core strategies across the primary
unflagged contamination scenario (\texttt{missed\_flags}, 1,000 cohorts,
where recording errors occur in unflagged records) and the mixed
gross-and-small error stress scenario
(\texttt{error\_gross\_and\_small}, 150 cohorts, combining \(10\times\)
gross transcription outliers with mild recording noise). The seven
strategies are: \textbf{Oracle} (regressing on true \(T_i\) for all
\(N = 600\) rows), \textbf{Flag Masking} (\texttt{v02\_mask\_normal},
setting flagged observations \(S_i = 1\) to \texttt{NA} for normal MI
while trusting unflagged observations \(S_i = 0\) as exact truth),
\textbf{Complete-Case} (\texttt{gold\_only}, analyzing only the
\(n_{\text{audit}} = 150\) audited rows), \textbf{Discard Recorded MI}
(\texttt{reference\_normal},
\texttt{catmech\_reference(use\_recorded\ =\ FALSE)}), \textbf{Simple
Linear MI} (\texttt{cole\_mi},
\texttt{catmech\_reference(use\_recorded\ =\ TRUE)}),
\textbf{Generalized Raking} (\texttt{raking}, two-phase calibration
estimator), and \textbf{Contamination Mixture}
(\texttt{catmech\_contamination()}):

{\def\LTcaptype{none} % do not increment counter
\begin{longtable}[]{@{}
  >{\raggedright\arraybackslash}p{(\linewidth - 14\tabcolsep) * \real{0.2812}}
  >{\centering\arraybackslash}p{(\linewidth - 14\tabcolsep) * \real{0.0833}}
  >{\centering\arraybackslash}p{(\linewidth - 14\tabcolsep) * \real{0.2188}}
  >{\centering\arraybackslash}p{(\linewidth - 14\tabcolsep) * \real{0.0833}}
  >{\centering\arraybackslash}p{(\linewidth - 14\tabcolsep) * \real{0.0833}}
  >{\centering\arraybackslash}p{(\linewidth - 14\tabcolsep) * \real{0.0833}}
  >{\centering\arraybackslash}p{(\linewidth - 14\tabcolsep) * \real{0.0833}}
  >{\centering\arraybackslash}p{(\linewidth - 14\tabcolsep) * \real{0.0833}}@{}}
\toprule\noalign{}
\begin{minipage}[b]{\linewidth}\raggedright
Scenario \& Strategy
\end{minipage} & \begin{minipage}[b]{\linewidth}\centering
\(\hat{\beta}\)
\end{minipage} & \begin{minipage}[b]{\linewidth}\centering
Bias (MCSE)
\end{minipage} & \begin{minipage}[b]{\linewidth}\centering
Emp SD
\end{minipage} & \begin{minipage}[b]{\linewidth}\centering
Mean SE
\end{minipage} & \begin{minipage}[b]{\linewidth}\centering
Cov.
\end{minipage} & \begin{minipage}[b]{\linewidth}\centering
Width
\end{minipage} & \begin{minipage}[b]{\linewidth}\centering
RMSE
\end{minipage} \\
\midrule\noalign{}
\endhead
\bottomrule\noalign{}
\endlastfoot
\textbf{Unflagged Errors (\(N_{\text{sim}} = 1,000\))} & & & & & & & \\
Oracle (Latent \(T\) Known) & 0.7978 & -0.0022 (0.0012) & 0.0394 &
0.0411 & 95.7\% & 0.1614 & 0.0394 \\
Flag Masking (\texttt{v02\_mask\_normal}) & 0.7185 & -0.0815 (0.0015) &
0.0472 & 0.0444 & 55.9\% & 0.1750 & 0.0942 \\
Complete-Case (\texttt{gold\_only}) & 0.7988 & -0.0012 (0.0026) & 0.0815
& 0.0837 & 95.7\% & 0.3309 & 0.0814 \\
Discard Recorded MI (\texttt{reference\_normal}) & 0.7896 & -0.0104
(0.0021) & 0.0676 & 0.0684 & 95.4\% & 0.2769 & 0.0684 \\
Simple Linear MI (\texttt{cole\_mi}) & 0.7806 & -0.0194 (0.0018) &
0.0557 & 0.0528 & 92.4\% & 0.2101 & 0.0590 \\
Generalized Raking (\texttt{raking}) & 0.7853 & -0.0147 (0.0020) &
0.0631 & 0.0589 & 92.9\% & 0.2331 & 0.0648 \\
\textbf{Contamination Mixture} & \textbf{0.7952} & \textbf{-0.0048
(0.0014)} & \textbf{0.0446} & \textbf{0.0467} & \textbf{94.9\%} &
\textbf{0.1842} & \textbf{0.0449} \\
\textbf{Gross \& Small (\(N_{\text{sim}} = 150\))} & & & & & & & \\
Oracle (Latent \(T\) Known) & 0.7970 & -0.0030 (0.0034) & 0.0411 &
0.0412 & 94.7\% & 0.1617 & 0.0411 \\
Flag Masking (\texttt{v02\_mask\_normal}) & 0.0609 & -0.7391 (0.0055) &
0.0674 & 0.0155 & 0.7\% & 0.0614 & 0.7422 \\
Complete-Case (\texttt{gold\_only}) & 0.7962 & -0.0038 (0.0064) & 0.0787
& 0.0845 & 95.3\% & 0.3339 & 0.0785 \\
Discard Recorded MI (\texttt{reference\_normal}) & 0.7890 & -0.0110
(0.0055) & 0.0671 & 0.0695 & 95.3\% & 0.2816 & 0.0678 \\
Simple Linear MI (\texttt{cole\_mi}) & 0.5646 & -0.2354 (0.0261) &
0.3200 & 0.0617 & 64.7\% & 0.2562 & 0.3964 \\
Generalized Raking (\texttt{raking}) & 0.5156 & -0.2844 (0.0281) &
0.3436 & 0.0693 & 53.7\% & 0.2740 & 0.4451 \\
\textbf{Contamination Mixture} & \textbf{0.7883} & \textbf{-0.0117
(0.0035)} & \textbf{0.0432} & \textbf{0.0436} & \textbf{94.0\%} &
\textbf{0.1714} & \textbf{0.0447} \\
\end{longtable}
}

In the \texttt{missed\_flags} scenario, trusting unflagged recordings
(\texttt{v02\_mask\_normal}) suffers severe attenuation bias
(\(-0.0815\)), collapsing 95\% coverage to 55.9\%. Complete-case
analysis (\texttt{gold\_only}) and discarding recorded values
(\texttt{reference\_normal}) remain unbiased but forfeit precision (mean
CI widths of 0.3309 and 0.2769). The \texttt{catmech} contamination
mixture model recovers unbiased estimates (bias: \(-0.0048\)), achieves
94.9\% coverage, and narrows confidence intervals by 33.5\% relative to
discarding recorded values (\texttt{reference\_normal}) and by 44.3\%
relative to complete-case analysis (\texttt{gold\_only}). When gross
transcription errors are mixed with small discrepancies
(\texttt{error\_gross\_and\_small}), both simple linear MI (bias:
\(-0.2354\), coverage: 64.7\%) and generalized raking (bias:
\(-0.2844\), coverage: 53.7\%) break down because gross outliers corrupt
the linear calibration slope, whereas \texttt{catmech\_contamination()}
preserves 94.0\% coverage and low RMSE (0.0447).

Across the four primary 1,000-cohort cells (\texttt{clean} with no
recording errors, \texttt{flagged\_errors} where errors are confined to
flagged rows, \texttt{missed\_flags} where errors occur in unflagged
rows, and \texttt{site\_shift\_declared} where site-level additive
shifts are combined with occasional larger errors),
\texttt{catmech\_contamination()} met the pre-specified
\([0.925, 0.975]\) Wilson coverage screen in three cells (96.5\%,
96.3\%, and 94.9\% coverage) and narrowly missed the screen in
\texttt{site\_shift\_declared} at 93.7\% coverage (Wilson interval
\([0.9202, 0.9505]\), bias \(-0.0085\), MCSE 0.0016), while reducing
mean interval width by 30\% to 38\% across all four primary cells
relative to \texttt{reference\_normal}.

To provide a balanced counterweight to the primary cells, the 23 stress
cells (150 cohorts each) deliberately violated model assumptions under a
pre-specified failure threshold of coverage \(< 90.0\%\) or
\(|\text{bias}| > 0.030\). Under this rule,
\texttt{catmech\_contamination()} failed in three stress cells:

\begin{enumerate}
\def\labelenumi{\arabic{enumi}.}
\tightlist
\item
  \textbf{Heavy-tailed target distribution
  (\texttt{heavy\_tail\_target})}: Bias reached \(+0.0330\) with 87.3\%
  coverage because the Gaussian mixture model misclassifies genuine
  heavy-tailed extremes in \(T_i\) as recording errors and shrinks them
  toward \(X_i \beta\); by contrast, \texttt{reference\_normal} (93.3\%
  coverage, bias \(+0.0130\)), \texttt{cole\_mi} (92.0\% coverage, bias
  \(-0.0003\)), and \texttt{raking} (92.0\% coverage, bias \(-0.0060\))
  held nominal coverage.
\item
  \textbf{Non-ignorable audit selection on latent truth with errors
  present (\texttt{nonignorable\_missed})}: When audit selection
  depended directly on the unobserved true value \(T_i\) in the presence
  of unflagged recording errors, \texttt{catmech\_contamination()}
  suffered bias of \(-0.0506\) and 79.3\% coverage (whereas
  design-consistent \texttt{raking} achieved bias \(+0.0070\) and 87.3\%
  coverage). When records were clean (\texttt{nonignorable\_clean}),
  coverage held at 92.5\% but 17 of 150 fits (11.3\%) were refused for
  slow MCMC mixing.
\item
  \textbf{Undeclared reference measurement noise
  (\texttt{noisy\_undeclared})}: When reference measurements \(R_i\)
  contained undeclared assay noise (\(\sigma_{\text{ref}} > 0\) left at
  the default \texttt{reference\_sd\ =\ 0}), bias was \(-0.0261\) and
  coverage fell to 88.7\% (triggering the near-universal-disagreement
  warning in all 150 cohorts). Declaring \texttt{reference\_sd}
  (\texttt{noisy\_declared}) restored bias to \(-0.0071\) and coverage
  to 91.8\% (147/150 returned, vs.~52.7\% coverage for
  \texttt{reference\_normal} and 38.0\% for \texttt{cole\_mi}).
\end{enumerate}

Two additional stress regimes were borderline or exhibited elevated
refusal rates: bimodal two-sided errors (\texttt{error\_bimodal}) was
right at the boundary with bias \(-0.0293\) and 90.0\% coverage;
range-check flags computed directly from recorded values
(\texttt{flag\_from\_recorded}) had 10 of 150 (6.7\%) convergence
refusals (94.3\% coverage among returned fits); and sparse 5\% reference
audits (\texttt{reference\_05pct}) had 17 of 150 (11.3\%) refusals (14
unsupported-stratum design refusals and 3 convergence refusals, with
91.7\% coverage among returned intervals). Finally, in an 8-scenario
evaluation of \texttt{catmech\_holdout\_check()} (40 cohorts per cell),
false alarms were absent under the true model (0/40 in \texttt{clean}
and \texttt{missed\_flags}), but alarm power against structural
misspecification was low: it flagged 18 of 40 cohorts (45\%) in
\texttt{site\_shift\_undeclared}, only 1 of 40 in
\texttt{heavy\_tail\_target} (with 6 unresolved fold failures), and 0 of
40 in \texttt{nonlinear\_target}, because held-out rows test
observation-level prediction rather than regression slope bias.

\subsection{Scenario 2: Finite-sample variance and clustered multilevel
imputation}\label{scenario-2-finite-sample-variance-and-clustered-multilevel-imputation}

In multicenter trials with \(C\) clinical sites, facility-level
clustering induces correlation that distorts naive variance estimators.
We evaluated 148,500 requests (148,486 returned intervals and 14
refusals: 12 rank-deficient audited designs on \texttt{iron} and 2
zero-residual-variance fits on \texttt{stillbirths}) across simulated
clustered cohorts and facility-level recounts of 11 maternal and child
health indicators from the Rwanda Health Management Information System
(HMIS) \citep{nshimyiryo2020health}. Standard small-sample
cluster-robust sandwich standard errors (\texttt{cr1s}) and
Bell-McCaffrey / Pustejovsky-Tipton bias-reduced linearization with
Satterthwaite degrees of freedom (\texttt{cr2}, which inflated interval
widths by up to 95\% under random-intercept imputation and passed only
23 of 44 simulated cells) failed the pre-specified eligibility rules.
The delete-one-cluster jackknife (\texttt{cv3}) restored 100\%
complete-data recount coverage (\texttt{gold\_only}) to 92.0\%--97.5\%
across the 11 Rwanda HMIS indicators and narrowly passed pre-specified
eligibility under random-intercept outcome imputation
(\texttt{catmech\_reference(cluster\ =\ )}, achieving \(\ge 92.0\%\)
coverage on 9 of 11 Rwanda indicators, with \texttt{syphilis} at 90.0\%
and \texttt{iron} at 91.9\%, and passing 36 of 44 simulated cells while
reducing mean distance from nominal 95\% coverage from 3.18 to 1.47
percentage points).

Specifically, this 148,500-request study also established two major
operational boundaries:

\begin{enumerate}
\def\labelenumi{\arabic{enumi}.}
\tightlist
\item
  \textbf{Imputation shortfall under linear MI on clustered recounts}:
  Applying \texttt{cv3} variance did not rescue unclustered linear MI
  (\texttt{catmech\_reference(use\_recorded\ =\ TRUE)} without
  \texttt{cluster}) on the 11 Rwanda HMIS indicators, where coverage
  remained below 92\% on 7 of 11 indicators (dropping to 81.0\% on
  \texttt{iron}, 81.6\% on \texttt{anc1\_early}, 84.2\% on
  \texttt{syphilis}, and 86.1\% on \texttt{pnc}), leaving the overall
  accuracy verdict as \texttt{PARTIAL} (and simulated rare-count
  recounts averaged 87.5\% coverage under \texttt{cv3}).
\item
  \textbf{Predictor-slope attenuation under clustered random-intercept
  MI}: When the audited variable \(T_i\) is a \textbf{predictor} in the
  outcome regression
  (\texttt{Y\ \textasciitilde{}\ T\ +\ x\ +\ stratum}) rather than the
  outcome, clustered random-intercept imputation
  (\texttt{catmech\_reference(cluster\ =\ )}) attenuated the estimated
  slope toward zero by up to \(-16.1\%\) with 84.4\% coverage at a 25\%
  audit under \texttt{cv3} (and \(-10.1\%\) bias, 86.0\% coverage at a
  50\% audit), whereas fitting
  \texttt{catmech\_reference(use\_recorded\ =\ TRUE)} without
  \texttt{cluster} and pooling with
  \texttt{catmech\_pool\_lm(cluster\ =\ )} yielded \(-3.0\%\) bias and
  95.4\% coverage. Accordingly,
  \texttt{catmech\_reference(cluster\ =\ )} remains experimental and
  should only be used when the audited variable is the analysis outcome
  (\texttt{catmech\_pool\_lm()} issues an explicit warning if
  \texttt{catmech\_reference(cluster\ =\ )} is used on an imputed
  predictor).
\end{enumerate}

\section{Real-world empirical
benchmarks}\label{real-world-empirical-benchmarks}

\subsection{Public paired-discrepancy registry and survey
benchmarks}\label{public-paired-discrepancy-registry-and-survey-benchmarks}

Because no public record-level clinical trial chart SDV dataset exists,
we first evaluated \texttt{catmech} on two large-scale public datasets
containing real paired recording discrepancies on every row without
synthetic injection (\texttt{validation/REAL\_DATA\_SYNTHESIS.md}):
ClinicalTrials.gov (\(N = 27,208\) studies, comparing sponsor-reported
log enrollment against participant-flow log Started counts in a
regression of log study duration in months on log enrollment, phase, and
sponsor class) and the National Health and Nutrition Examination Survey
(NHANES) 2017--2018 \citep{nhanes2018} (\(N = 3,908\) adults, comparing
self-reported weight against examination-measured weight in kg in a
regression of systolic blood pressure in mmHg on weight, age, sex, and
race). Across 33,800 planned requests (1,000 resampled cohorts of size
\(N = 600\) at the primary 25\% audit fraction
\(n_{\text{audit}} = 150\), plus 200 cohorts per cell at 5\%, 10\%,
50\%, and targeted audit fractions), 33,613 intervals were returned
alongside 76 unsupported-design refusals, 83 MCMC convergence refusals,
and 28 singular raking errors:

{\def\LTcaptype{none} % do not increment counter
\begin{longtable}[]{@{}
  >{\raggedright\arraybackslash}p{(\linewidth - 12\tabcolsep) * \real{0.2232}}
  >{\centering\arraybackslash}p{(\linewidth - 12\tabcolsep) * \real{0.1161}}
  >{\centering\arraybackslash}p{(\linewidth - 12\tabcolsep) * \real{0.1339}}
  >{\centering\arraybackslash}p{(\linewidth - 12\tabcolsep) * \real{0.1161}}
  >{\centering\arraybackslash}p{(\linewidth - 12\tabcolsep) * \real{0.1696}}
  >{\centering\arraybackslash}p{(\linewidth - 12\tabcolsep) * \real{0.0982}}
  >{\centering\arraybackslash}p{(\linewidth - 12\tabcolsep) * \real{0.1429}}@{}}
\toprule\noalign{}
\begin{minipage}[b]{\linewidth}\raggedright
Dataset
\end{minipage} & \begin{minipage}[b]{\linewidth}\centering
Metric
\end{minipage} & \begin{minipage}[b]{\linewidth}\centering
Naive (Trust Rec.)
\end{minipage} & \begin{minipage}[b]{\linewidth}\centering
Discard Rec.
\end{minipage} & \begin{minipage}[b]{\linewidth}\centering
Contamination
\end{minipage} & \begin{minipage}[b]{\linewidth}\centering
Linear MI
\end{minipage} & \begin{minipage}[b]{\linewidth}\centering
Full Ref. (Oracle)
\end{minipage} \\
\midrule\noalign{}
\endhead
\bottomrule\noalign{}
\endlastfoot
\textbf{ClinicalTrials.gov} (\(N = 27,208\)) & Bias & +0.0043 & -0.0003
& +0.0068 & +0.0001 & -0.0003 \\
Discrepancy rate: 14.2\% & Coverage & 91.4\% & 92.3\% & 91.2\% & 91.8\%
& 91.2\% \\
Target: 0.0357 log mo / log N & Mean Width & 0.097 & 0.205 & 0.106 &
0.102 & 0.096 \\
& RMSE & 0.0282 & 0.0539 & 0.0308 & 0.0286 & 0.0275 \\
\textbf{NHANES 2017--2018} (\(N = 3,908\)) & Bias & -0.0180 & -0.0063 &
-0.0066 & -0.0025 & +0.0001 \\
Discrepancy rate (\(>1\) kg): 67.3\% & Coverage & 92.7\% & 94.2\% &
94.5\% & 93.9\% & 94.9\% \\
Target: 0.127 mmHg / kg & Mean Width & 0.149 & 0.297 & 0.155 & 0.154 &
0.140 \\
& RMSE & 0.0422 & 0.0785 & 0.0408 & 0.0409 & 0.0357 \\
\end{longtable}
}

On both datasets, retaining recorded values reduced RMSE by 43\% to 48\%
and narrowed confidence intervals by 48\% compared to discarding
unaudited values (\texttt{Discard\ Rec.}, \texttt{reference\_normal}).
Note that on ClinicalTrials.gov, the Full-Reference Oracle itself
covered 91.2\% because homoskedastic OLS standard errors (0.0244)
slightly underestimated empirical sampling variability (0.0275) under
non-constant residual variance, a gap inherited by all methods.
Moreover, on ClinicalTrials.gov at a 25\% audit, naively trusting
recorded values (\texttt{Naive\ (Trust\ Rec.)}) achieved an RMSE of
0.0282 and bias of \(+0.0043\), whereas
\texttt{catmech\_contamination()} had slightly higher bias (\(+0.0068\))
and RMSE (\(0.0308\)), and at a sparse 5\% audit fraction
\texttt{catmech\_contamination()} was weakest (returning 164 of 200
intervals {[}18\% refusal rate{]} with bias \(+0.0159\) and RMSE
\(0.0423\), vs.~200/200, bias \(+0.0062\), and RMSE \(0.0356\) for
\texttt{Linear\ MI}). On NHANES at 25\%, where self-reported weight bias
was substantial (\(-0.0180\), or 14\% of the target slope), both
\texttt{Linear\ MI} (bias \(-0.0025\), RMSE \(0.0409\)) and
\texttt{Contamination} (bias \(-0.0066\), RMSE \(0.0408\)) corrected the
naive bias (and at a 10\% NHANES audit, \texttt{Contamination} achieved
lower RMSE {[}\(0.0435\) vs.~\(0.0490\){]} than \texttt{Linear\ MI} with
14/200 refusals).

Most importantly, because real-world discrepancies in both datasets were
small relative to the conditional spread of \(T_i \mid X_i\),
\textbf{per-record posterior corruption probabilities \(p_i\) had
near-chance discrimination on real data}: AUC for identifying discrepant
unaudited rows at a 25\% audit was \textbf{0.543} on ClinicalTrials.gov
(Brier score 0.127 vs.~0.121 for base-rate prevalence alone; posterior
mean RMSE 0.277 vs.~0.252 for raw recorded values) and \textbf{0.500} on
NHANES (Brier score 0.222 vs.~0.220 for prevalence; posterior mean RMSE
4.81 vs.~4.73 kg for raw recorded values), compared to an AUC of 0.962
in the synthetic \texttt{missed\_flags} cell. Consequently, individual
record-level error probabilities \(p_i\) must not be used to adjudicate
individual patient records on diffuse real-world data. While the KS
alarm in \texttt{catmech\_holdout\_check()} stayed quiet on these
cohorts (firing in only 2 of 38 completed checks on ClinicalTrials.gov
and 1 of 40 on NHANES), the continuous held-out Brier and RMSE skill
scores correctly indicated that the mixture model had no predictive edge
over raw recorded values.

\subsection{Real-audit regression slope recovery
benchmarks}\label{real-audit-regression-slope-recovery-benchmarks}

We further conducted a pre-specified 45,600-request real-audit benchmark
(\texttt{validation/REAL\_AUDIT\_SYNTHESIS.md}) evaluating regression
slope recovery on the Rwanda HMIS facility register recounts
\citep{nshimyiryo2020health} alongside 498 paired effect-size recounts
from psychology meta-analyses \citep{maassen2020reproducibility}. Under
three pre-specified decision rules at a 25\% audit fraction:

\begin{enumerate}
\def\labelenumi{\arabic{enumi}.}
\tightlist
\item
  \textbf{Rule 1 (Does a partial 25\% audit with \texttt{Linear\ MI}
  beat raw recorded data?)}: A 25\% audit with
  \texttt{catmech\_reference(use\_recorded\ =\ TRUE)} paid clearly on
  only \textbf{1 of 4 Rwanda slopes} (\texttt{ANC\ 4th\ visits}, where
  errors scaled with facility size {[}\(r = +0.477\){]} and recovery
  RMSE fell to \(0.458\times\) that of raw recorded data), tied on
  \texttt{PNC} (\(0.923\times\) in the initial 500 cohorts,
  \(1.009\times\) {[}\(0.969, 1.048\){]} on 2,000 fresh seeds, and
  \(1.136\times\) against the full-pool slope), and increased recovery
  RMSE relative to raw recorded values on \texttt{Syphilis}
  (\(2.393\times\)) and \texttt{Iron} (\(1.125\times\)).
\item
  \textbf{Rule 2 (Does \texttt{catmech\_contamination()} add value over
  \texttt{Linear\ MI}?)}: The pre-specified verdict on Rwanda HMIS was
  \textbf{\texttt{INCONCLUSIVE}}. Although
  \texttt{catmech\_contamination()} had lower point recovery RMSE than
  \texttt{Linear\ MI} on \texttt{PNC} (\(0.053\) vs.~\(0.072\), ratio
  \(0.745\)) and \texttt{Iron} (\(0.112\) vs.~\(0.156\), ratio
  \(0.717\)), it failed the pre-specified coverage condition on
  \textbf{all three primary slopes} and failed the \(\ge 90\%\) interval
  return-rate condition on \textbf{two of three}: clustered 95\%
  coverage was \textbf{71.6\%} on \texttt{PNC} (468/500 {[}93.6\%{]}
  returned, vs.~85.4\% coverage for \texttt{Linear\ MI} and 95.2\% for
  the recount oracle), \textbf{66.6\%} on \texttt{Syphilis} (404/500
  {[}80.8\%{]} returned, RMSE \(0.195\) vs.~\(0.182\) for
  \texttt{Linear\ MI} and \(0.076\) for raw recorded), \textbf{35.4\%}
  on \texttt{Iron} (424/500 {[}84.8\%{]} returned, vs.~79.5\% coverage
  for \texttt{Linear\ MI}), and \textbf{49.7\%} on the size-dependent
  \texttt{ANC\ 4th\ visit} stress slope (RMSE \(0.124\) vs.~\(0.085\)
  for \texttt{Linear\ MI}, retaining 63\% of reported bias). Across all
  20 Rwanda cells, \texttt{catmech\_contamination()} was closest to 95\%
  coverage in only 1 of 20 cells (worst case 35.4\%), compared to 10 of
  20 cells for \texttt{Linear\ MI} (worst case 79.5\%), and per-record
  error AUCs were 0.501 to 0.582.
\item
  \textbf{Rule 3 (Psychology effect-size recount controls where reported
  bias is near zero)}: On the 498 psychology effect-size recounts
  \citep{maassen2020reproducibility}, where raw reported slopes already
  closely matched recounted truth, \texttt{Linear\ MI} at a 25\% audit
  inflated recovery RMSE relative to raw recorded values by
  \textbf{\(2.705\times\)} (standardized mean difference) and
  \textbf{\(1.578\times\)} (correlation \(r\)), whereas
  \texttt{catmech\_contamination()} stayed close to the recorded values
  (\(1.306\times\) and \(0.950\times\)).
\end{enumerate}

Together, these real-audit results demonstrate that the advantage of
\texttt{catmech\_contamination()} under gross, separable errors remains
\textbf{synthetic-only} (because Rwanda's recount discrepancies are
moderate and non-separable), and that on moderate real-world audit
discrepancies \texttt{Linear\ MI} provides substantially more reliable
interval coverage than the contamination mixture model.

\section{Methodological decision
framework}\label{methodological-decision-framework}

Based on our synthetic and empirical findings, we recommend the
following evidence-based protocol for clinical data audits:

\begin{verbatim}
         [ Clinical Measurement Discrepancy Detected ]
                              |
   1. Screen error rate & stratum heterogeneity: catmech_recording_audit()
   2. Compare candidate strategies on audited rows: catmech_compare_strategies()
                              |
         Are discrepancies gross, large, and separable
     (with Gaussian target residuals & ignorable audit selection)?
                         /           \
                       YES            NO (diffuse noise, rounding,
                       /                  moderate recount diffs,
        [ Contamination Mixture ]         or heavy-tailed target)
         catmech_contamination()                   \
         (verify >=90% return rate        [ Simple Linear MI ]
          & holdout Brier/RMSE skill)      catmech_reference(
                                             use_recorded = TRUE)
\end{verbatim}

\begin{enumerate}
\def\labelenumi{\arabic{enumi}.}
\tightlist
\item
  \textbf{Gross, separable contamination}: When recording errors
  represent large, discrete corruption events (e.g., decimal
  transcription slips or \(10\times\) unit errors) and target residuals
  are approximately Gaussian under ignorable audit selection,
  \texttt{catmech\_contamination()} prevents severe attenuation bias
  while maintaining near-oracle precision.
\item
  \textbf{Subtle, diffuse, or moderate real-world discrepancies}: When
  errors represent continuous rounding, mild reporting bias, or moderate
  non-separable recount differences (as in ClinicalTrials.gov, NHANES,
  and Rwanda HMIS), simple linear imputation
  (\texttt{catmech\_reference(use\_recorded\ =\ TRUE)}) is the
  recommended baseline, avoiding MCMC convergence refusals and mixture
  interval undercoverage. When the audited variable is a predictor in a
  clustered study, fit
  \texttt{catmech\_reference(use\_recorded\ =\ TRUE)} without
  \texttt{cluster} and pool with
  \texttt{catmech\_pool\_lm(cluster\ =\ )}.
\end{enumerate}

\section{Discussion: simulation-to-practice divergence and
methodological
implications}\label{discussion-simulation-to-practice-divergence-and-methodological-implications}

A central contribution of this work is documenting the divergence
between synthetic Monte Carlo simulations and empirical real-world
audits, including full disclosure of pre-specified negative and
inconclusive findings.

In the statistical methodology literature, novel measurement-error and
imputation models are predominantly benchmarked on synthetic datasets
where contamination is gross, bimodal, and generated from the exact
parametric family assumed by the estimator. Under those synthetic
conditions (\texttt{missed\_flags} and
\texttt{error\_gross\_and\_small}), \texttt{catmech\_contamination()}
eliminates the severe attenuation bias that compromises simple linear
imputation (bias of \(-0.0117\) vs.~\(-0.2354\) in
\texttt{error\_gross\_and\_small}) while narrowing confidence intervals
by 44.3\% relative to complete-case analysis in \texttt{missed\_flags}
(\(0.1842\) vs.~\(0.3309\)). Even in synthetic stress testing, however,
our pre-specified counterfactual scenarios expose clear failure
boundaries: heavy-tailed target distributions
(\texttt{heavy\_tail\_target}, 87.3\% coverage), audit selection on the
latent true value (\texttt{nonignorable\_missed}, 79.3\% coverage), and
undeclared reference noise (\texttt{noisy\_undeclared}, 88.7\%
coverage).

Our empirical investigations across ClinicalTrials.gov, NHANES
2017--2018, and the Rwanda HMIS facility recount audit reinforce a
sobering practical lesson: real-world recording discrepancies are rarely
gross or separable into discrete clean and contaminated components. On
ClinicalTrials.gov and NHANES, simple linear MI
(\texttt{catmech\_reference(use\_recorded\ =\ TRUE)}) matched or
slightly beat the contamination mixture model in RMSE (0.0286 vs.~0.0308
on ClinicalTrials.gov; 0.0409 vs.~0.0408 on NHANES) with zero MCMC
overhead, while on Rwanda HMIS the contamination mixture model suffered
severe clustered interval undercoverage (35.4\% to 71.6\%) and 6\% to
19\% convergence refusal rates. Across all three real datasets,
per-record corruption probabilities \(p_i\) performed near chance (AUC
0.500 to 0.582).

This simulation-to-practice divergence demonstrates why \texttt{catmech}
is designed around an empirical diagnostic workflow rather than a single
flagship estimator:

\begin{enumerate}
\def\labelenumi{\arabic{enumi}.}
\tightlist
\item
  Exploratory discrepancy screening via
  \texttt{catmech\_recording\_audit()} to quantify error magnitudes and
  between-stratum heterogeneity before choosing an imputation strategy;
\item
  Out-of-sample predictive cross-validation via
  \texttt{catmech\_holdout\_check()}, comparing held-out Brier and RMSE
  skill scores against raw recorded values rather than relying on the
  low-power KS uniformity alarm alone;
\item
  Direct side-by-side strategy comparison via
  \texttt{catmech\_compare\_strategies()} so study teams can verify
  whether the data support mixture modeling, linear MI, complete-case
  analysis, or unadjusted recorded values.
\end{enumerate}

\section{Regulatory documentation and data
governance}\label{regulatory-documentation-and-data-governance}

Clinical trial analyses are expected to be documented and reproducible:
FDA 21 CFR Part 11 \citep{fda2003part11} governs electronic records and
audit trails, and ICH E9 \citep{ich1998e9} asks that data handling and
statistical methods be prespecified and documented.
\texttt{catmech\_audit\_trail()} supports that documentation. It does
not certify compliance, which depends on the validated system, access
controls, and procedures around the analysis. It generates provenance
records in Markdown, JSON, or standalone HTML, capturing:

\begin{enumerate}
\def\labelenumi{\arabic{enumi}.}
\tightlist
\item
  Full execution environment metadata (R version, operating system
  platform, package version, UTC timestamp).
\item
  SHA-256 hashes of the input data and of each of the \(m\) completed
  datasets, computed over the serialized R data frames (so they can be
  recomputed only in R, not with a file-hashing tool), and of the model
  specification as canonical JSON text. A rerun can be compared with
  these hashes; they do not cover the pooled results, and they do not
  make the record tamper-evident, because anyone who edits the data can
  also regenerate the record.
\item
  MCMC convergence diagnostics recorded on the audit trail
  (\texttt{convergence\_passed} pass/fail status and
  \texttt{mcmc\_iterations} count, with full Gelman-Rubin \(\hat{R}\)
  and bulk/tail effective sample size vectors stored in
  \texttt{result\$model\_diagnostics}).
\item
  Pooled parameter estimates and multi-scenario sensitivity envelopes.
\end{enumerate}

\section{Quality assurance and
verification}\label{quality-assurance-and-verification}

\texttt{catmech} has been developed under an automated software
verification harness:

\begin{itemize}
\tightlist
\item
  \textbf{Unit tests}: 199 \texttt{testthat} tests with 1,812
  expectations, achieving 0 failures, 0 warnings, and 0 skips.
\item
  \textbf{CRAN conformance}: \texttt{R\ CMD\ check\ -\/-as-cran} reports
  0 errors and 0 warnings.
\item
  \textbf{Automated artifact verification}: Verified against an
  independent Python verification suite (\texttt{verify-artifacts.py},
  \texttt{verify-release-archive.py}, and
  \texttt{verify\_paper\_numbers.py}) confirming exact correspondence
  across 10 frozen simulation and empirical benchmark archives and over
  500,000 replicate requests.
\end{itemize}

\section{Availability and software
details}\label{availability-and-software-details}

\texttt{catmech} is open-source software distributed under the MIT
license. It requires R (\(\ge\) 4.1.0) and a C++17 compiler, utilizing
the non-base CRAN packages \texttt{Rcpp}, \texttt{RcppArmadillo},
\texttt{digest}, \texttt{jsonlite}, \texttt{mice}, \texttt{pan},
\texttt{posterior}, \texttt{survival}, \texttt{truncnorm}, and
\texttt{yaml} (alongside base R packages \texttt{graphics},
\texttt{grDevices}, \texttt{parallel}, \texttt{stats}, and
\texttt{utils}). Source code, documentation, and vignettes are available
at \url{https://github.com/patrickiben/catmech}.

\renewcommand\refname{References}
\bibliography{paper.bib}

\end{document}
